\documentclass[11pt]{article}
\usepackage[T1]{fontenc}
\usepackage{lmodern}
\usepackage[margin=1in]{geometry}
\usepackage{amsmath,amssymb,amsthm,mathtools,booktabs,array}
\usepackage[hidelinks]{hyperref}
\newtheorem{theorem}{Theorem}[section]
\newtheorem{corollary}[theorem]{Corollary}
\newtheorem{lemma}[theorem]{Lemma}
\newtheorem{proposition}[theorem]{Proposition}
\theoremstyle{remark}
\newtheorem{remark}[theorem]{Remark}
\newcommand{\R}{\mathbb R}
\newcommand{\wind}{\operatorname{wind}}

\title{A positive-Jacobian noninjective sigmoid ridge map\\and the sharp hidden-width threshold}
\author{Logan M. Dixon\thanks{Independent Researcher. Email: \texttt{lmdixon23@gmail.com}. ORCID: 0009-0001-0592-462X.}}
\date{September 2026}

\begin{document}
\maketitle

\begin{abstract}
The Neural Jacobian Conjecture asks whether an everywhere-positive Jacobian determinant forces global injectivity for affine-ridge logistic maps from $\R^n$ to $\R^n$. We give an exact rational counterexample at the smallest pair left open in that formulation, $n=2$, $N=4$. Its determinant is positive everywhere by a two-bridge arithmetic--geometric-exponential decomposition based on exact affine relations among the logits. Two disjoint interval Krawczyk certificates establish noninjectivity directly. Independently, its row-ordered asymptotic octagon has winding number two at an explicit rational target; the companion planar degree--winding theorem then shows that this target has exactly two preimages. A block construction lifts the example to every $n\ge2$ at width $N=n+2$, and hidden-unit duplication gives every larger width. Combined with the companion injectivity theorem at widths $N=n$ and $N=n+1$, this gives the sharp hidden-width threshold. We also prove a closed-form two-bridge determinant criterion, an explicit multiparameter counterexample family with nonzero bias freedom, and an extension to every saturating activation whose derivative is an even log-concave probability density strictly positive on all of $\R$. Finally, an exact computer-assisted classification on the open refined magnitude chambers of the enumerated canonical direction realizations shows that nonnegative winding spectrum with multiplicity at least two can occur only with the alternating determinant-sign pattern realized by the counterexample.
\end{abstract}

\medskip
\noindent\textbf{2020 Mathematics Subject Classification.} 26B10 (primary); 55M25, 65G20, 68T07, 15A15.

\smallskip
\noindent\textbf{Keywords.} global injectivity; sigmoid ridge network; Neural Jacobian Conjecture; Krawczyk operator; winding number; validated numerics.

\section{Main results}
The Neural Jacobian Conjecture was formulated by Chen and Jiang~\cite{ChenJiang}. Their work proves the case $N=n+1$ and leaves the higher-width regime $N\ge n+2$ unresolved. Let $\sigma(t)=(1+e^{-t})^{-1}$ act componentwise.

\begin{theorem}[Planar four-ridge counterexample]\label{thm:main}
Set
\[
 A=\begin{pmatrix}
 1&0&-8&2\\
 0&1&-1&1/8
 \end{pmatrix},\qquad
 B=\begin{pmatrix}
 1&0\\
 0&1\\
 -1/5&5/4\\
 -1/2&9/5
 \end{pmatrix}.
\]
Then $F(x)=A\sigma(Bx)$ satisfies
\[
 \det DF(x)>0\qquad(x\in\R^2),
\]
but $F$ is not injective. More precisely, for
\[
 y=\left(-\frac{11}{2},\frac{1}{10}\right)
\]
one has $\#F^{-1}(y)=2$.
\end{theorem}

\begin{corollary}[Sharp hidden-width threshold]\label{cor:sharp}
For every $n\ge2$ and every $N\ge n+2$, there is a width-$N$ affine-ridge logistic map $G:\R^n\to\R^n$ with $\det DG>0$ everywhere that is not injective. For $N<n$, a nonsingular Jacobian is impossible by rank, while at $N=n$ and $N=n+1$ every such map is injective~\cite[Thm.~3.2]{DixonPolygon}. Thus $N=n+2$ is the sharp onset of counterexamples. The two closed widths are settled by a graph-coordinate reduction in which the determinant is affine in a single varying slot; the present counterexample shows that the resulting injectivity conclusion fails at $N=n+2$. For $n=1$, positivity of the derivative implies injectivity at every width.
\end{corollary}

The proof separates three logically distinct claims. Sections~\ref{sec:age}--\ref{sec:det} prove pointwise determinant positivity. Proposition~\ref{prop:krawczyk} certifies two distinct preimages directly, so the conjecture refutation does not depend on the companion paper. The companion asymptotic-polygon results~\cite{DixonPolygon} are used to upgrade those two preimages to exact global multiplicity two and, through the closed-width theorem at $N=n$ and $N=n+1$, to obtain sharpness.

\section{Bridge AGE inequalities}\label{sec:age}
The determinant argument applies to a broad derivative class.

\begin{lemma}[Bridge AGE inequality with bias]\label{lem:age}
Let $\rho(t)=C e^{-\psi(t)}$, where $C>0$ and $\psi:\R\to\R$ is even, convex, nonnegative, and satisfies $\psi(0)=0$. Suppose
\[
 z_k=\alpha z_i+\beta z_j+\delta,
 \qquad
 a=|\alpha|,\quad b=|\beta|,\quad r=1-a-b\ge0.
\]
Choose $w_i,w_j>0$ with
\[
 w_i+w_j=1,\qquad w_i\ge a,\qquad w_j\ge b.
\]
If $r=0$, additionally assume $\delta=0$. Define
\[
 E=\begin{cases}
 \exp\!\left[-r\psi\!\left(\delta/r\right)\right],&r>0,\\
 1,&r=0.
 \end{cases}
\]
If $P,Q,M>0$ and
\[
 \left(\frac{P}{w_j}\right)^{w_j}
 \left(\frac{Q}{w_i}\right)^{w_i}E>M,
\]
then
\[
 P\rho(z_i)\rho(z_k)+Q\rho(z_j)\rho(z_k)-M\rho(z_i)\rho(z_j)>0.
\]
\end{lemma}

\begin{proof}
After division by $\rho(z_i)\rho(z_j)>0$, weighted AM--GM gives
\[
 P\frac{\rho(z_k)}{\rho(z_j)}
 +Q\frac{\rho(z_k)}{\rho(z_i)}
 \ge
 \left(\frac{P}{w_j}\right)^{w_j}
 \left(\frac{Q}{w_i}\right)^{w_i}
 \frac{\rho(z_k)}{\rho(z_i)^{w_i}\rho(z_j)^{w_j}}.
\]
Evenness and convexity give
\[
 \psi(z_k)\le a\psi(z_i)+b\psi(z_j)+r\psi(\delta/r)
\]
when $r>0$, with the evident zero-slack form when $r=0$. Since $w_i\ge a$, $w_j\ge b$, and $\psi\ge0$, the derivative ratio is at least $E$. The strict coefficient inequality completes the proof.
\end{proof}

\begin{theorem}[Closed-form two-bridge criterion]\label{thm:two-bridge}
Suppose a planar four-ridge determinant has the alternating form
\[
\begin{aligned}
J={}&p_{01}\rho_0\rho_1-n_{02}\rho_0\rho_2+p_{03}\rho_0\rho_3\\
&+p_{12}\rho_1\rho_2-n_{13}\rho_1\rho_3+p_{23}\rho_2\rho_3,
\end{aligned}
\]
where every displayed magnitude is positive. Suppose
\[
 z_1=\alpha_0z_0+\alpha_2z_2+\delta_1,
 \qquad
 z_2=\beta_1z_1+\beta_3z_3+\delta_2,
\]
and choose admissible weights $u_0+u_2=1$ and $v_1+v_3=1$ as in Lemma~\ref{lem:age}. Let $E_1,E_2$ be the corresponding bias penalties and put
\[
 L_1=u_0\left[
 \frac{n_{02}}
 {E_1(p_{01}/u_2)^{u_2}}
 \right]^{1/u_0},
\]
\[
 L_2=v_3\left[
 \frac{n_{13}}
 {E_2(p_{23}/v_1)^{v_1}}
 \right]^{1/v_3}.
\]
If
\[
 p_{12}>L_1+L_2,
\]
then $J>0$ everywhere.
\end{theorem}

\begin{proof}
Choose $s$ with $L_1<s<p_{12}-L_2$. Apply Lemma~\ref{lem:age} to
\[
 p_{01}\rho_0\rho_1+s\rho_1\rho_2-n_{02}\rho_0\rho_2
\]
and to
\[
 (p_{12}-s)\rho_1\rho_2+p_{23}\rho_2\rho_3-n_{13}\rho_1\rho_3.
\]
Both groups are strictly positive, and the unused term $p_{03}\rho_0\rho_3$ is positive.
\end{proof}

\begin{remark}
Theorem~\ref{thm:two-bridge} is a two-negative-term AGE certificate in the sense of relative-entropy and SAGE nonnegativity methods~\cite{ChandrasekaranShah,MurrayEtAl}. Here the row circuits reduce the certificate to explicit scalar inequalities and a one-dimensional allocation.
\end{remark}

\section{Everywhere-positive determinant}\label{sec:det}
Write the columns of $A$ as $a_0,\ldots,a_3$ and the rows of $B$ as $b_0,\ldots,b_3$. In pair order
\[
01,02,03,12,13,23,
\]
the output minors are
\[
 1,-1,\frac18,8,-2,1,
\]
and the row minors are
\[
 1,\frac54,\frac95,\frac15,\frac12,\frac{53}{200}.
\]
Thus the Cauchy--Binet coefficients are
\[
 1,-\frac54,\frac9{40},\frac85,-1,\frac{53}{200}.
\]
Put $z=Bx$ and $\rho_i=\sigma'(z_i)$. Since
\[
 \sigma'(t)=\frac14e^{-\psi(t)},
 \qquad \psi(t)=2\log\cosh(t/2),
\]
Lemma~\ref{lem:age} applies. Split
\[
 \frac85=\frac38+\frac{49}{40}
\]
and write
\[
 \det DF=G_1+G_2+\frac9{40}\rho_0\rho_3,
\]
where
\[
 G_1=\rho_0\rho_1+\frac38\rho_1\rho_2-\frac54\rho_0\rho_2
\]
and
\[
 G_2=\frac{49}{40}\rho_1\rho_2
 +\frac{53}{200}\rho_2\rho_3-\rho_1\rho_3.
\]
The rows give the exact relations
\[
 z_1=\frac4{25}z_0+\frac45z_2,
 \qquad
 z_2=\frac{53}{100}z_1+\frac25z_3.
\]
For $G_1$, use weights $w_0=1/5$, $w_2=4/5$. The coefficient bound is
\[
 \left(\frac{15}{8}\right)^{1/5}
 \left(\frac54\right)^{4/5}>\frac54.
\]
For $G_2$, use weights $w_1=53/100$, $w_3=47/100$. The required inequality is
\[
 \left(\frac12\right)^{53/100}
 \left(\frac{245}{94}\right)^{47/100}>1,
\]
which is equivalent to
\[
 245^{47}>94^{47}2^{53}.
\]
Hence $G_1,G_2>0$, proving $\det DF>0$ everywhere.

\section{Exact winding and two preimages}\label{sec:winding}
The row walls are pairwise distinct. Over one semicircle their crossing order gives the signed edge sequence
\[
 +a_1,+a_2,+a_3,-a_0,-a_1,-a_2,-a_3,+a_0.
\]
Starting at $a_0$, the asymptotic octagon has vertices
\[
 (1,0),(1,1),(-7,0),(-5,1/8),
\]
\[
 (-6,1/8),(-6,-7/8),(2,1/8),(0,0).
\]
The target $y=(-11/2,1/10)$ is not on an edge. Exact rational ray crossing gives
\[
 \wind(\gamma,y)=2.
\]
The following validated certificate already proves noninjectivity, without using the companion degree--winding theorem.

\begin{proposition}[Two certified preimages]\label{prop:krawczyk}
Let
\[
\begin{aligned}
   x_1&=(-0.5652157847725153,\ 3.1135129207599814),\\
   x_2&=(25.1754764056295704,\ 7.4516846575961514),
 \end{aligned}
\]
and let $X_i$ be the closed interval box of radius $10^{-10}$ about $x_i$ in each coordinate. Then $X_1$ and $X_2$ are disjoint and each contains exactly one solution of $F(x)=y$. Consequently $F$ is not injective.
\end{proposition}

\begin{proof}
Put $G(x)=F(x)-y$. For a box $X$ with midpoint $m$ and a nonsingular preconditioner $C$, the Krawczyk operator is
\[
 K_G(X)=m-C\,G(m)+\bigl(I-C\,DG(X)\bigr)(X-m).
\]
The standard Krawczyk inclusion test $K_G(X)\subseteq\operatorname{int}X$ certifies a unique zero of $G$ in $X$~\cite[Ch.~8]{MooreInterval}. The maintained certificate evaluates the inclusion with outward-rounded Arb ball arithmetic at 256-bit precision. Using the exact terminating decimals displayed above as the box centres, the two infinity-norm residuals are below $3.1\times10^{-17}$, the centre Jacobian determinants are approximately $5.5264\times10^{-3}$ and $6.4604\times10^{-4}$, and the Krawczyk images lie strictly inside the corresponding boxes coordinatewise. The boxes are separated by more than $25$ in the first coordinate, hence are disjoint.
\end{proof}

\begin{remark}
This certificate uses only the matrices of Theorem~\ref{thm:main}, the target $y$, and outward-rounded interval arithmetic. It therefore establishes noninjectivity independently of the asymptotic geometry. The winding computation supplies the exact global multiplicity.
\end{remark}

The companion degree--winding theorem states, for continuously saturating planar ridge maps with pairwise nonparallel rows and nonvanishing Jacobian,
\[
 \#F^{-1}(y)=\operatorname{sgn}(\det DF)\wind(\gamma,y)
\]
for $y$ off the asymptotic loop~\cite{DixonPolygon}. All hypotheses have been established, and the orientation is positive, so $\#F^{-1}(y)=2$. Thus the two certified roots are the only roots of $F(x)=y$, completing Theorem~\ref{thm:main}.

\section{An explicit multiparameter counterexample family}\label{sec:family}
The counterexample belongs to a transparent region rather than merely a small certified ball.

\begin{theorem}[Explicit analytic family]\label{thm:family}
Let
\[
 A(u,t)=\begin{pmatrix}
 1&0&-8&u\\
 0&1&-1&t
 \end{pmatrix}
\]
and keep $B$ as in Theorem~\ref{thm:main}. For $c=(c_0,c_1,c_2,c_3)$ define
\[
 \delta_1=c_1-\frac4{25}c_0-\frac45c_2,
 \qquad
 \delta_2=c_2-\frac{53}{100}c_1-\frac25c_3.
\]
If
\[
 \frac{199}{100}<u<\frac{51}{25},
 \qquad
 \frac3{25}<t<\frac{13}{100},
\]
and
\[
 |\delta_1|<\frac3{40},
 \qquad
 |\delta_2|<\frac1{20},
\]
then
\[
 F_{u,t,c}(x)=A(u,t)\sigma(Bx+c)
\]
has positive Jacobian determinant everywhere and winding number two at $y=(-11/2,1/10)$. Consequently it is not injective and $\#F_{u,t,c}^{-1}(y)=2$.
\end{theorem}

\begin{proof}
The mixed coefficients are
\[
 1,-\frac54,\frac95t,\frac85,-\frac u2,
 \frac{53}{200}(u-8t).
\]
Use the same split $8/5=3/8+49/40$. For the first bridge, the zero-bias coefficient ratio divided by $5/4$ is
\[
 R_1=\left(\frac32\right)^{1/5}.
\]
The exact inequality
\[
 3\cdot37^5>2\cdot40^5
\]
gives $R_1>40/37$, while $e^{3/40}<40/37$ follows from $e^x<(1-x)^{-1}$ for $0<x<1$. Since the logistic potential satisfies $\psi(s)\le |s|$, the first bridge remains strict under $|\delta_1|<3/40$.

For the second bridge, the zero-bias ratio is
\[
 R_2(u,t)=
 \frac{(245/94)^{47/100}((u-8t)/2)^{53/100}}{u/2}.
\]
On the stated rectangle,
\[
 \frac{\partial}{\partial u}\log R_2
 =\frac{800t-47u}{100u(u-8t)}>0,
 \qquad
 \frac{\partial}{\partial t}\log R_2<0.
\]
Thus the infimum on the open rectangle is the value at the closure corner $(u,t)=(199/100,13/100)$. Raising to the hundredth power gives an exact integer comparison showing this lower bound exceeds $20/19$. Since $e^{1/20}<20/19$, the second bridge remains strict under $|\delta_2|<1/20$.

The asymptotic polygon is independent of $c$. Target-edge incidence in the moving $(u,t)$ family occurs, in the relevant component, on the lines
\[
 t=\frac1{10},\qquad u=15t,\qquad u=\frac52.
\]
The parameter rectangle lies strictly in the component containing $(2,1/8)$, so the target winding remains two. The companion degree--winding formula~\cite{DixonPolygon} then gives exactly two preimages and finishes the proof.
\end{proof}

\begin{remark}
Theorem~\ref{thm:family} gives open freedom in two output parameters and all four biases subject to two open linear residual conditions. The general criterion in Theorem~\ref{thm:two-bridge} gives a still larger implicit analytic class in the full network parameters.
\end{remark}

\section{Sharp width classification}\label{sec:lift}
For $n\ge2$, define
\[
 G_n(x_1,\ldots,x_n)
 =\bigl(F(x_1,x_2),\sigma(x_3),\ldots,\sigma(x_n)\bigr).
\]
This is a one-hidden-layer affine-ridge network of width $4+(n-2)=n+2$. Its Jacobian is block diagonal, so
\[
 \det DG_n(x)=\det DF(x_1,x_2)\prod_{k=3}^n\sigma'(x_k)>0.
\]
It is noninjective because the two planar preimages can be augmented with identical remaining coordinates. For $N>n+2$, repeatedly duplicate a hidden unit with nonzero output column $a$ and replace that column by two copies $a/2,a/2$ attached to identical rows and biases. The represented function and Jacobian are unchanged. The duplicated rows are parallel, so the standing nondegeneracy hypothesis of the polygon and degree results no longer holds; neither is reapplied after duplication, and positivity and noninjectivity transfer because the function itself is unchanged. This proves Corollary~\ref{cor:sharp}.

\section{Symmetric log-concave activation class}\label{sec:activation}
The logistic formula is not essential.

\begin{corollary}[Activation-class extension]\label{cor:activation}
Let $q(t)=C e^{-\psi(t)}$ be an even log-concave probability density that is strictly positive on all of $\R$, so that $\psi:\R\to\R$ is finite valued, normalized so that $\psi\ge0$ and $\psi(0)=0$, and define
\[
 f(t)=\int_{-\infty}^{t}q(s)\,ds.
\]
Then the matrices in Theorem~\ref{thm:main}, with $\sigma$ replaced by $f$, define an everywhere-positive-Jacobian noninjective ridge map with two preimages of the same target. The dimension and width lift also applies.
\end{corollary}

\begin{proof}
The derivative $f'=q$ satisfies Lemma~\ref{lem:age}. The asymptotic polygon and its winding use only continuity, strict monotonicity, and saturation to zero and one. The determinant argument is therefore unchanged, and the companion degree--winding theorem~\cite{DixonPolygon} applies with the same asymptotic loop and target, giving exactly two preimages.
\end{proof}

This includes the logistic sigmoid, $(1+\tanh t)/2$, the Gaussian distribution function, and the normalized generalized Gaussian family with density proportional to $e^{-|t|^p}$ for $p\ge1$.

\section{Alternating-sign chamber classification}\label{sec:classification}
The explicit counterexample is sufficient to refute the conjecture. The following exact classification places its two negative pairs within the canonical direction realizations enumerated by the certificate.

Fix an ordered generic four-generator centrally symmetric octagon with edge sequence
\[
 g_0,g_1,g_2,g_3,-g_0,-g_1,-g_2,-g_3.
\]
Write $g_i=\lambda_i d_i$ with fixed pairwise nonparallel directions and $\lambda_i>0$. Here generic magnitudes mean that $\lambda$ lies off every nonzero $L$ and $K$ wall in Proposition~\ref{prop:chambers}; the classification below is restricted to these open refined chambers.

\begin{proposition}[Refined magnitude chambers]\label{prop:chambers}
For fixed directed edge directions, there is a finite central hyperplane arrangement in positive magnitude space such that, on each open chamber, the crossing set, the order of crossings along every edge, polygon simplicity, and the winding spectrum are constant.
\end{proposition}

\begin{proof}
Every vertex is homogeneous linear in $\lambda$. Edge--vertex orientation predicates are signs of homogeneous linear forms $L_{rs}(\lambda)$. If two nonparallel comparison edges meet the line of edge $r$ at parameters $t_{rs}$ and $t_{rt}$, the numerator of $t_{rs}-t_{rt}$ is another homogeneous linear form $K_{r;s,t}(\lambda)$. Add all nonzero $L$ and $K$ walls. Within an open chamber, no vertex crosses a nonincident edge and no intersection points exchange order along an edge. The induced planar immersion is isotopic, so its face winding data are constant.
\end{proof}

\begin{theorem}[Canonical-realization alternating classification]\label{thm:classification}
Fix the canonical direction realization of each labeled uniform rank-two chirotope used by the certificate below, and restrict $\lambda$ to an open chamber of its refined magnitude arrangement. In the chosen edge order and positive orientation, suppose the winding spectrum is nonnegative and contains a value at least two. Then the six generator determinant signs, in pair order
\[
 01,02,03,12,13,23,
\]
are
\[
 (+,-,+,+,-,+).
\]
\end{theorem}

\begin{proof}[Computer-assisted exact proof]
The certificate enumerates all $48$ labeled uniform realizable rank-two chirotopes on four generators. For each chirotope it constructs every $L$ and $K$ wall from Proposition~\ref{prop:chambers}. Exact Gordan-circuit alternatives certify feasibility or infeasibility of each chamber sign vector. Each feasible chamber receives a rational interior point, and exact segment and ray-crossing predicates compute its winding spectrum. Across all chirotopes there are $2448$ feasible refined chambers. Exactly one sign pattern admits a nonnegative spectrum with maximum at least two, namely the displayed alternating pattern. For that pattern, $140$ chambers have spectrum $\{0,1,2\}$, $51$ have spectrum $\{0,1,2,3\}$, and $14$ have spectrum $\{-1,0,1\}$.
\end{proof}

\begin{remark}
The classification is not load-bearing for Theorem~\ref{thm:main}. Within their open refined magnitude chambers, the enumerated canonical realizations identify the only determinant-sign pattern that can support nonnegative winding multiplicity at least two. The quantifier does not extend to every generic realization: four generic directions in $\mathbb{RP}^1$ retain a one-dimensional modulus modulo $\mathrm{PGL}(2,\R)$, since that group is sharply $3$-transitive, and the winding spectrum is not a chirotope invariant --- positive rescaling can preserve every pairwise determinant sign while changing a simple loop into a self-intersecting one~\cite[Prop.~7.5]{DixonPolygon}. Extending the statement to all realizations would require adding that modulus to the exact decomposition. The full rational certificate and replay are supplied separately.
\end{remark}

\section{Verification and reproducibility}
The accompanying repository contains exact rational checks of the minors, Pl\"ucker relation, row circuits, coefficient inequalities, wall order, polygon vertices, target exclusion, winding number, the multiparameter family, the dimension lift, and the finite chamber classification. The two preimage boxes are certified by outward-rounded Arb arithmetic. A separate internal SymPy reconstruction rebuilds the central matrices, row relations, wall order, vertices, and winding with a separately written implementation; it is corroborating evidence rather than an external review.

The finite chamber theorem is computer assisted. Its completeness claim depends on the exact Gordan-circuit enumeration and is logically separate from the counterexample. The conjecture refutation itself uses the analytic determinant proof together with Proposition~\ref{prop:krawczyk}. The companion degree--winding theorem is needed for the stronger exact-multiplicity statement, and the companion closed-width theorem is additionally needed for the sharp-threshold corollary. Numerical root finding is retained only as a diagnostic and is not used as proof evidence.

\section{Conclusion}
An everywhere-positive Jacobian does not force injectivity for sigmoid ridge networks once the hidden width reaches $n+2$. The obstruction already occurs at $n=2$, $N=4$: a two-bridge analytic certificate keeps the local orientation positive while two disjoint Krawczyk boxes rigorously certify a collision. The companion degree--winding theorem shows that the same target has exactly two preimages, and its closed-width theorem turns the construction into a sharp threshold separating $N\le n+1$ from $N\ge n+2$ for $n\ge2$. The mechanism persists on a transparent multiparameter region and across all saturating activations whose derivative is an even log-concave probability density strictly positive on all of $\R$. On the open refined magnitude chambers of the enumerated canonical direction realizations, the classification singles out the alternating two-negative determinant-sign pattern as the only one supporting a nonnegative winding spectrum with multiplicity at least two.

\begin{thebibliography}{99}
\bibitem{ChenJiang}
X. Chen and X. Jiang,
\emph{Moonshine: An Autonomous Mathematical Research Agent Centered on Conjecture Generation},
arXiv:2606.10806, 2026.

\bibitem{DixonPolygon}
L. M. Dixon,
\emph{Asymptotic-polygon global inversion for planar saturating ridge networks and an injective neural--Jacobian separation witness},
preprint, 2026.

\bibitem{MooreInterval}
R. E. Moore, R. B. Kearfott, and M. J. Cloud,
\emph{Introduction to Interval Analysis},
SIAM, Philadelphia, 2009.

\bibitem{ChandrasekaranShah}
V. Chandrasekaran and P. Shah,
\emph{Relative Entropy Relaxations for Signomial Optimization},
SIAM Journal on Optimization 26 (2016), 1147--1173.

\bibitem{MurrayEtAl}
R. Murray, V. Chandrasekaran, and A. Wierman,
\emph{Newton Polytopes and Relative Entropy Optimization},
Foundations of Computational Mathematics 21 (2021), 1703--1737.

\end{thebibliography}

\end{document}
